\documentclass[letterpaper, onecolumn, draftclsnofoot,10pt, compsoc]{IEEEtran}
\usepackage{graphicx}
\usepackage[hyphens]{url}
\usepackage{setspace}

%\usepackage{biblatex}
%\addbibresource{references.bib}

\usepackage{geometry}
\geometry{textheight=9.5in, textwidth=7in}

% 1. Fill in these details
\def \CapstoneTeamName{		The Polycules}
\def \CapstoneTeamNumber{		63}
\def \GroupMemberOne{			Joshua Lioy}
\def \GroupMemberTwo{			Corynna Park}
\def \GroupMemberThree{			Jackson Wells}
\def \CapstoneProjectName{		Molecules in 3D?! And in Color!? That I can Hold in My Hand? No way!!!}
\def \CapstoneSponsorCompany{		Oregon State University}
\def \CapstoneSponsorPerson{		Victor Hsu}

% 2. Uncomment the appropriate line below so that the document type works
\def \DocType{			%Problem Statement
				%Requirements Document
				Technology Review
				%Design Document
				%Progress Report
				}
			
\newcommand{\NameSigPair}[1]{\par
\makebox[2.75in][r]{#1} \hfil 	\makebox[3.25in]{\makebox[2.25in]{\hrulefill} \hfill		\makebox[.75in]{\hrulefill}}
\par\vspace{-12pt} \textit{\tiny\noindent
\makebox[2.75in]{} \hfil		\makebox[3.25in]{\makebox[2.25in][r]{Signature} \hfill	\makebox[.75in][r]{Date}}}}
% 3. If the document is not to be signed, uncomment the RENEWcommand below
%\renewcommand{\NameSigPair}[1]{#1}

%%%%%%%%%%%%%%%%%%%%%%%%%%%%%%%%%%%%%%%
\begin{document}
\begin{titlepage}
    \pagenumbering{gobble}
    \begin{singlespace}
    	\includegraphics[height=4cm]{coe_v_spot1}
        \hfill 
        % 4. If you have a logo, use this includegraphics command to put it on the coversheet.
        %\includegraphics[height=4cm]{CompanyLogo}   
        \par\vspace{.2in}
        \centering
        \scshape{
            \huge CS Capstone \DocType \par
            {\large\today}\par
            \vspace{.5in}
            \textbf{\Huge\CapstoneProjectName}\par
            \vfill
            {\large Prepared for}\par
            \Huge \CapstoneSponsorCompany\par
            \vspace{5pt}
            {\Large\NameSigPair{\CapstoneSponsorPerson}\par}
            {\large Prepared by }\par
            Group\CapstoneTeamNumber\par
            % 5. comment out the line below this one if you do not wish to name your team
            \CapstoneTeamName\par 
            \vspace{5pt}
            {\Large
               % \NameSigPair{\GroupMemberOne}\par
                \NameSigPair{\GroupMemberTwo}\par
               % \NameSigPair{\GroupMemberThree}\par
            }
            \vspace{20pt}
        }
        \begin{abstract}
        % 6. Fill in your abstract    
        Many different technologies are involved in the process of 3D printing, including the type of filament used to print, the file type used to store data, and the slicing software used to communicate instructions to the printer.
     	This document will review the aforementioned three pieces of technology and how they will affect this project. 
        Through giving three different options for each piece of technology, using explicit criteria to evaluate, and analyzing each option using outside research, the best-suited options for the project will be determined.   
        \end{abstract}     
    \end{singlespace}
\end{titlepage}
\newpage
\pagenumbering{arabic}
\tableofcontents
% 7. uncomment this (if applicable). Consider adding a page break.
%\listoffigures
%\listoftables
\clearpage

% 8. now you write!
\section{Project Overview} 
Currently, the capability to produce monochromatic 3D models already exists. 
However, not much is known about how to print polychromatic 3D models.
Multicolor 3D printing will enhance learning and foster better absorption of concepts by creating better visual representations of data. 
In order to complete this project, a robust workflow will need to be developed which will allow anyone with basic knowledge of the 3D printing process to print multicolor models.
To ensure availability and accessibility, all the software used will be free and open-source.
Several pieces of technology will be involved in this project and influence the outcome. 
The pieces that will be discussed in this document are the filament type, file type, and slicing software.


\section{3D Filament Type}
\subsection{Overview}
3D printing filament is the thermoplastic feedstock used in the 3D printing process.
The type of filament used directly affects the physical properties of the 3D object that is printed. 
Currently, there are many different filament types that can be used for 3D printing, each with different attributes and different temperature requirements.
In this section, three filament types, PLA, ABS, and Nylon 618, will be compared to determine the type that is best suited for this project. 
\subsection{Criteria}
When analyzing the different filament types, the criteria that will be looked at will be: 
\begin{itemize}
\item Strength
\item Flexibility
\item Durability
\item Difficulty to Print
\item Print Temperature
\item Bed Temperature
\item Other Notable Qualities
\end{itemize}
\subsection{Potential Choices}
\subsubsection{Choice One: Polylactic Acid}
The first type of filament we will be addressing is polylactic acid, or PLA. 
PLA is a commonly used material for 3D printing. 
It is biodegradable and derived from renewable resources; thus, it is more environmentally friendly than most plastic materials. 
Considering the criteria, PLA is of medium strength, has low flexibility, and medium durability. 
PLA is relatively easy to print and requires a print temperature range between 180-220 degrees Celsius.  
When using PLA, a heated bed is not needed. 
With some filament types, heated beds are needed to improve print quality and prevent warping that deforms the model.\cite{heatbedweb}
A notable quality about PLA is its biocompatibility with a human body, meaning it does not produce a toxic or immunological response when exposed to the body or bodily fluids.
Because of this, PLA is safe to use and could be utilized for medically-related 3D printing. \cite{filamentweb}
\subsubsection{Choice Two: Acrylonitrile Butadiene Styrene}
When 3D printers first came around, acrylonitrile butadiene styrene, or ABS, was the most popularly used material; it is still commonly used today. 
In terms of the criteria, ABS is of medium strength, but is weaker than PLA.
ABS has medium flexibility and high durability.
It requires a print temperature of about 210-250 degrees Celsius.
Because of its higher temperature requirement, ABS is of a medium level in terms of printing difficulty. 
Unlike PLA, ABS requires a heated bed with a temperature of about 50-100 degrees Celsius.
Something to note is that when printing with ABS, fumes are produced which can be harmful for people or pets with breathing difficulties.
Thus, when ABS is in use, one should be printing in a well-ventilated area and avoid inhaling the fumes. 
Another quality to note is that when considering the cost of 3D materials, ABS is the cheapest filament type. \cite{filamentweb}
\subsubsection{Choice Three: Nylon 618}
The third filament type that will be considered is Nylon, specifically Nylon 618. 
Compared to PLA and ABS, Nylon materials are much less common and harder to find, but its printing attributes are comparable.
Nylon 618 is of high strength and durability, due to its self-bonding properties. 
The material itself comes in a natural white color. 
It has high flexibility and requires a print temperature of 220-260 degrees Celsius, giving it a medium print difficulty. 
A heated bed is not required when using Nylon 618.
A notable quality of Nylon 618 is that the material is able to be dyed any color desirable using simple fabric dye. 
This feature can be utilized when wanting to print multicolor objects. 
Another quality to note is that Nylon materials absorb moisture from air, thus, when not being used they should be kept in a dry environment. \cite{filamentweb}

\subsection{Discussion}
Since this project intends to be useful to a wide audience, some of which may be limited in terms of the financing and time they have available to dedicate to 3D printing, the filament type best suited for it would be one that is easily accessible, affordable, and easy to print.
For the project's polychromatic purposes, a filament that can be represent different colors in a 3D model conveniently and quickly is desirable.
Compared to the other filaments, Nylon 618 is less accessible, more costly, and relatively difficult to print.
In addition, to represent different colors, Nylon 618 must be dyed; this would add a significant time requirement to an already time-intensive process.
While ordering Nylon that has already been dyed would seem like a solution, due to its lack of accessibility and popularity, it is unlikely that this is a feasible option.
In comparing the remaining two filament types, PLA and ABS, PLA is easier to print due to its lower temperature requirements. 
PLA also does not require a heated bed, unlike ABS. 
Another considerable advantage of PLA over ABS is its biocompatibility; PLA is safer to use and does not produce harmful fumes.
\subsection{Conclusion}
After comparing the advantages and disadvantages of ABS, PLA, and Nylon 618, the filament type that should be used to complete this project is PLA.

\section{3D File Type}
\subsection{Overview}
Before an object can be printed in 3D, its information must first be encoded and stored in a file. 
Depending on the file type, this information is saved as either plain text or binary data and represents physical details about an object such as is shape, coloring, texture, shadings, etc. \cite{all3dpweb}
There are many different file types exported by CAD (Computer-Aided Drawing) software to be used in 3D printing, and these types vary in the amount of information they store.
The following section will analyze and compare three file types, Stereolithography, Object, and 3D Manufacturing Format, to determine which will be the best to use for this project.
 
\subsection{Criteria}
When analyzing the different file types, the criteria that will be looked at will be: 
\begin{itemize}
\item Complexity
\item Representation
\item Difficulty to Export
\item Format Type
\end{itemize}
\subsection{Potential Choices}
\subsubsection{Choice One: Stereolithography}
Stereolithography, or STL, is one of the oldest and most popularly used 3D file formats in 3D printing. 
Files of this type are indicated with an .STL extension.
STL is limited in that it only stores information about a 3D object's geometry.
Considering the complexity criteria, the STL file format is one of the simplest 3D file formats available.
In terms of representation, it only contains information about an object's geometry, ignoring other physical attributes, including color.
Since STL files store less information, they are relatively small.
STL files are supported by most 3D printers and easily exported by most CAD software; the export feature for STL is usually set by default. \cite{3dbeginweb}
The STL file format type is neutral, meaning files in this format are open source and can be used by any software.
\cite{all3dpweb}
\subsubsection{Choice Two: Object File}
OBJ, or object file, is another commonly used file type when it comes to 3D printing. 
These types of files are indicated with the .OBJ extension and specify both binary and ASCII encodings of data.
In terms of complexity, OBJ files are considered to be more complex than STL files because they can store a greater amount of information.
Like STL files, fies in the OBJ format can store information on a 3D object's geometry. 
In addition, they are also capable of including data about the object's appearance, such as color and texture. 
OBJ files are easily exported by most CAD software and supported by 3D printers; however, some require a plug-in. \cite{3dbeginweb}
The file format type for OBJ depends on its variant; its ASCII variant is neutral, but its binary variant is proprietary. 
A file of proprietary type has readability that is restricted to only its own and other permitted software. 
The ASCII variant is able to be read and used by any software.
\cite{all3dpweb}

\subsubsection{Choice Three: 3D Manufacturing Format}
The third file type we will take into consideration is 3MF, the 3D Manufacturing Format.
3MF is a new 3D printing format that is XML-based (Extensible Markup Language), human-readable, and allows design applications to export complete information on an object to a wide variety of applications, platforms, services, and printers. \cite{3mfweb}
This file type is simple with its clear and concise specification, but also complex because it stores a significant amount of data.
3MF files are able to fully describe a model, including object geometry, appearance, and more, while narrowing the scope of the data to avoid adding unnecessary complexity. 
Since 3MF is a newer file format type, not every software supports it as of yet.
However, there is an open source library as well as Microsoft implemented Windows 10 APIs which allow for the reading, writing, and validating of 3MF files. 
3MF was developed as a neutral format type, as one of the main objectives in its development was interoperability.
This means that 3MF files are fully open source and free to use. \cite{3mfweb}
\subsection{Discussion}
For this project's purposes, the file type that will be best-suited to solve the problem would be one that contains information about an object's geometry and appearance, as multicolor representation is our main goal.
The file type must also be supported by the slicing software utilized as well as free and open source.
Potentially, 3MF could have several advantages over OBJ and STL due to it being most recently developed.
According to Andrew Zaleski's article, "Why These Big Companies Want a New 3D File Format," the STL and OBJ file formats are "outdated and clunky" and may have "interoperability issues when used by some of the newer 3D printers available." \cite{newformatweb}
Since STL files do not support multicolor representation, despite being the simplest to use, STL as the file type would seem to be the least feasible option.
In order to use STL files to create a multicolor model, a method to splice together several STL files which represent each separately colored section would have to be developed. This may go beyond the time limitations and scope of the project. 
Moving forward, in comparing OBJ and 3MF files, they both are able to contain color information.
However, OBJ files may be more complex than 3MF files, as OBJ files are not designed for manufacturing and may contain unnecessary information. \cite{3mfweb}
Both may require additional steps when implementing their formats, and both have the necessary resources to do so.

\subsection{Conclusion}
After comparing the three file types, STL, OBJ, and 3MF, the one that should be utilized for this project is 3MF, as it appears to be the best for multicolor representation and interoperability. 

\section{3D Slicing Software}
\subsection{Overview}
In order for the 3D printer to understand and print an object, the file containing the data about the object must be translated.
This translation is done through slicing software.
Slicing software turns the information about an object's 3D attributes into G-code, a control language for machines. 
This code essentially creates instructions for the 3D printer to follow. \cite{whatisweb}
To determine a slicing software that will be suitable for this project, Voxelizer, Simplify3D, and Cura will be discussed in the subsections that follow. 

\subsection{Criteria}
When analyzing the different file types, the criteria that will be looked at will be: \begin{itemize}
\item Features
\item File Formats Supported
\item Compatibility with Operating Systems
\item Compatibility with 3D Printers
\item Cost

\end{itemize}
\subsection{Potential Choices}
\subsubsection{Choice One: Voxelizer 2}
Voxelizer 2 is the slicing software provided by the creators of the Zmorph 3D printer, which is the printer that will be used for this project. 
The features of the software, as listed on its website, include on-the-spot voxel editing, dedicated multi-material workflows, each with special presets, ambient occlusion for 3D printing which improves the object quality and aesthetic, innovative image mapping including color blending, and multiple file format support. 
It supports multimaterial and 3D printing, allowing up to 8 colors in a single print as well as filament mixing.
The file formats that is supports are STL, JPG, DXF, and DICOM medical files. 
Voxelizer is compatible with both Windows and Mac OS operating systems. 
Despite being the slicing software created for the Zmorph 3D printer, Voxelizer is able to be used with several other desktop 3D printers, including Prusa i3, BCN Sigma, and any RepRap machine. 
The software itself is free to the public and easy to download. \cite{voxelweb}

\subsubsection{Choice Two: Simplify3D}
The next slicing software that will be discussed is Simplify3D.
According to the website for the software, Simplify3D has "model setup and plating, slicing and print file creation, pre-print simulations, customizable support structures, mesh analysis and repair, machine control and monitoring, and much more."
Simplify3D supports STL, OBJ, and 3MF file formats and has versions to support Windows, Mac, and Linux operating systems.
Already tested and optimized for hundreds of different 3D printers, the Simplify3D website even includes a "Check Compatibility" feature to see if a specific 3D printer is supported.
Simplify3D does come at a cost of \$149 USD, however there are educational discounts for teachers, administrators, and students. \cite{3dsweb}

\subsubsection{Choice Three: Cura}
The Ultimaker Cure software is advertised as "the world's most advanced 3D printer software," suitable for both the novice and expert. 
Cura's special features, as listed on the Ultimaker website, include 200 customizable settings when optimizing 3D models, the ability to print multiple objects at once, each with their own settings, and downloadable plugins to "create seamless integration with leading design and engineering software."
The file formats supported are STL, 3MF, and OBJ.
Although designed to be combined with Ultimaker printers, Cura supports many other 3D printers and is customizable so that almost any 3D printer can be supported if it is not already.
Compatible with Mac OS, Windows, and Linux, Cura is completely free and open-source. \cite{curaweb}

\subsection{Discussion}
For this project, it is a requirement that all software involved must be free and open-source.
Since Simplify3D has a cost, it is not able to be used for this project. 
In the previous section, it was decided that 3MF will be the file format used throughout this project. 
Voxelizer does not support the 3MF file format, thus, it appears that Cura would be the best slicing software to use for this project. 

\subsection{Conclusion}
After taking Voxelizer, Simplify3D, and Cura into consideration, Cura will be the best slicing software to use for this project.

%\printbibliography
\bibliographystyle{IEEEtran}
\bibliography{references}{}


\end{document}
